\documentclass[10pt,a4paper]{amsart}
\usepackage{amsmath,amssymb,amsthm,graphicx}
% Retained minimal preamble for the Stacks Project chapter flat.tex (More on Flatness).
% Verbatim excerpt of the chapter preamble: Xy-pic setup (native preamble lines 14-18) and the native
% theorem declarations with the mathematical macros (native preamble lines 154-217). Copied mechanically;
% nothing added, removed or redefined. The wrapper class and the amsmath/amssymb/amsthm/graphicx base
% packages are supplied by the assembler, so the class/externaldocument part of the native preamble is omitted.
% For commutative diagrams we use Xy-pic
\usepackage[all]{xy}

% We use 2cell for 2-commutative diagrams.
\xyoption{2cell}
% Theorem environments.
%
\theoremstyle{plain}
\newtheorem{theorem}[subsection]{Theorem}
\newtheorem{proposition}[subsection]{Proposition}
\newtheorem{lemma}[subsection]{Lemma}

\theoremstyle{definition}
\newtheorem{definition}[subsection]{Definition}
\newtheorem{example}[subsection]{Example}
\newtheorem{exercise}[subsection]{Exercise}
\newtheorem{situation}[subsection]{Situation}

\theoremstyle{remark}
\newtheorem{remark}[subsection]{Remark}
\newtheorem{remarks}[subsection]{Remarks}

\numberwithin{equation}{subsection}

% Macros
%
\def\lim{\mathop{\mathrm{lim}}\nolimits}
\def\colim{\mathop{\mathrm{colim}}\nolimits}
\def\Spec{\mathop{\mathrm{Spec}}}
\def\Hom{\mathop{\mathrm{Hom}}\nolimits}
\def\Ext{\mathop{\mathrm{Ext}}\nolimits}
\def\SheafHom{\mathop{\mathcal{H}\!\mathit{om}}\nolimits}
\def\SheafExt{\mathop{\mathcal{E}\!\mathit{xt}}\nolimits}
\def\Sch{\mathit{Sch}}
\def\Mor{\mathop{\mathrm{Mor}}\nolimits}
\def\Ob{\mathop{\mathrm{Ob}}\nolimits}
\def\Sh{\mathop{\mathit{Sh}}\nolimits}
\def\NL{\mathop{N\!L}\nolimits}
\def\CH{\mathop{\mathrm{CH}}\nolimits}
\def\proetale{{pro\text{-}\acute{e}tale}}
\def\etale{{\acute{e}tale}}
\def\QCoh{\mathit{QCoh}}
\def\Ker{\mathop{\mathrm{Ker}}}
\def\Im{\mathop{\mathrm{Im}}}
\def\Coker{\mathop{\mathrm{Coker}}}
\def\Coim{\mathop{\mathrm{Coim}}}

% Boxtimes
%
\DeclareMathSymbol{\boxtimes}{\mathbin}{AMSa}{"02}

%
% Macros for moduli stacks/spaces
%
\def\QCohstack{\mathcal{QC}\!\mathit{oh}}
\def\Cohstack{\mathcal{C}\!\mathit{oh}}
\def\Spacesstack{\mathcal{S}\!\mathit{paces}}
\def\Quotfunctor{\mathrm{Quot}}
\def\Hilbfunctor{\mathrm{Hilb}}
\def\Curvesstack{\mathcal{C}\!\mathit{urves}}
\def\Polarizedstack{\mathcal{P}\!\mathit{olarized}}
\def\Complexesstack{\mathcal{C}\!\mathit{omplexes}}
% \Pic is the operator that assigns to X its picard group, usage \Pic(X)
% \Picardstack_{X/B} denotes the Picard stack of X over B
% \Picardfunctor_{X/B} denotes the Picard functor of X over B
\def\Pic{\mathop{\mathrm{Pic}}\nolimits}
\def\Picardstack{\mathcal{P}\!\mathit{ic}}
\def\Picardfunctor{\mathrm{Pic}}
\def\Deformationcategory{\mathcal{D}\!\mathit{ef}}

% --- appended 2026-10-12 (Chief): the project own definitions of the reference/slogan/history
% display environments, copied mechanically from the shared Stacks preamble lines 8-12
% (\usepackage{verbatim} plus the three \newenvironment lines, which the chapter preamble keeps
% verbatim from the project preamble). Nothing added, removed or redefined: these environments
% carry no mathematical content and the project itself typesets them as comments. Without these
% three lines a counted statement whose own text uses them fails with Environment slogan undefined.
\usepackage{verbatim}
\newenvironment{reference}{\comment}{\endcomment}
\newenvironment{slogan}{\comment}{\endcomment}
\newenvironment{history}{\comment}{\endcomment}
\usepackage{cleveref}
\crefname{originaldefinition}{original definition}{original definitions}
\Crefname{originaldefinition}{Original definition}{Original definitions}
\makeatletter
\newlabel{more-morphisms-definition-relatively-finitely-presented-sheaf}{{05H1}{source}{}{sourcepointer-1}{}}
\newlabel{more-morphisms-definition-relatively-finitely-presented-sheaf@cref}{{[originaldefinition][0][]05H1}{[1][0][]source}{}{}{}}
\makeatother
\crefname{originallemma}{original lemma}{original lemmas}
\Crefname{originallemma}{Original lemma}{Original lemmas}
\makeatletter
\newlabel{more-morphisms-lemma-finite-morphism-relative-finite-presentation}{{09T9}{source}{}{sourcepointer-2}{}}
\newlabel{more-morphisms-lemma-finite-morphism-relative-finite-presentation@cref}{{[originallemma][0][]09T9}{[1][0][]source}{}{}{}}
\makeatother
\crefname{originallemma}{original lemma}{original lemmas}
\Crefname{originallemma}{Original lemma}{Original lemmas}
\makeatletter
\newlabel{more-morphisms-lemma-relative-finite-presentation}{{09T8}{source}{}{sourcepointer-3}{}}
\newlabel{more-morphisms-lemma-relative-finite-presentation@cref}{{[originallemma][0][]09T8}{[1][0][]source}{}{}{}}
\makeatother
\crefname{originaltheorem}{original theorem}{original theorems}
\Crefname{originaltheorem}{Original theorem}{Original theorems}
\makeatletter
\newlabel{theorem-finite-type-flat}{{05IK}{source}{}{sourcepointer-4}{}}
\newlabel{theorem-finite-type-flat@cref}{{[originaltheorem][0][]05IK}{[1][0][]source}{}{}{}}
\makeatother
\title{Openness of the flat locus in the finite type case, via the finite presentation theorem}
\author{The Stacks Project Authors}
\date{stacks-project:flat.tex:local-snapshot-20261011; GNU FDL 1.2 or later}
\begin{document}
\maketitle
\noindent\textit{Editorial note.} All delivered mathematical text is the author's own, copied from the pinned Stacks Project chapter file flat.tex; no statement or proof step was rewritten, reordered or removed.\par\medskip
\noindent\textit{Editorial note.} Preamble provenance: the chapter's own retained minimal preamble (sources/stacks-flat/retained-preamble.tex, lines 1-85) is a verbatim excerpt of the native Xy-pic setup and of the native theorem declarations and mathematical macros, followed by the project's own reference/slogan/history display-environment definitions copied mechanically from the shared preamble; the wrapper class and the amsmath/amssymb/amsthm/graphicx base packages are supplied by the assembler, and extra\_packages is left empty because everything needed is already in the retained preamble. No macro was invented or redefined.\par\medskip
\noindent\textit{Editorial note.} Same-chapter pointers are located at their \textbackslash{}label lines inside the pinned chapter file; cross-chapter pointers are located in the project's own tag list staged in the snapshot (sources/stacks-chow/tags.txt), at the verbatim lines carrying their canonical tags.\par\medskip
\noindent\textit{Allowed standard prerequisites (editorial).}\begin{itemize}
\item The finite presentation version of openness of flatness. Same chapter flat.tex, label theorem-finite-type-flat, invoked at proof line 4250. Applied to the ambient affine space over the base and the pushforward of the module, it gives that the flat locus there is open and that the restriction of the pushforward is of finite presentation.
\item Relative finite presentation and its behaviour under a finite morphism. More on Morphisms, labels more-morphisms-definition-relatively-finitely-presented-sheaf (statement, proof line 4243), more-morphisms-lemma-relative-finite-presentation and more-morphisms-lemma-finite-morphism-relative-finite-presentation (proof lines 4266-4267). Supplies the notion of being locally finitely presented relative to the base and shows that it is inherited from the pushforward along the closed immersion.
\end{itemize}\medskip
\ifdefined\hypertarget\hypertarget{sourcepointer-1}{}\fi
\noindent\textit{Source reference.} Original definition 05H1: sources/stacks-chow/tags.txt, source lines 7080--7080. Navigation pointer for \textbackslash{}ref\{more-morphisms-definition-relatively-finitely-presented-sheaf\} used inside the counted span (line 4244); the argument that consumes it is retained in the span.\par\medskip
\ifdefined\hypertarget\hypertarget{sourcepointer-2}{}\fi
\noindent\textit{Source reference.} Original lemma 09T9: sources/stacks-chow/tags.txt, source lines 12382--12382. Navigation pointer for \textbackslash{}ref\{more-morphisms-lemma-finite-morphism-relative-finite-presentation\} used inside the counted span (line 4267); the argument that consumes it is retained in the span.\par\medskip
\ifdefined\hypertarget\hypertarget{sourcepointer-3}{}\fi
\noindent\textit{Source reference.} Original lemma 09T8: sources/stacks-chow/tags.txt, source lines 12381--12381. Navigation pointer for \textbackslash{}ref\{more-morphisms-lemma-relative-finite-presentation\} used inside the counted span (line 4266); the argument that consumes it is retained in the span.\par\medskip
\ifdefined\hypertarget\hypertarget{sourcepointer-4}{}\fi
\noindent\textit{Source reference.} Original theorem 05IK: sources/stacks-flat/pinned-original-source.tex, source lines 4097--4097. Navigation pointer for \textbackslash{}ref\{theorem-finite-type-flat\} used inside the counted span (line 4230); the argument that consumes it is retained in the span.\par\medskip
\begin{lemma}[Finite type version of Theorem \ref{theorem-finite-type-flat}]
\label{lemma-finite-type-flat}
Let $f : X \to S$ be a morphism of schemes.
Let $\mathcal{F}$ be a quasi-coherent $\mathcal{O}_X$-module.
Assume
\begin{enumerate}
\item $X \to S$ is locally of finite type,
\item $\mathcal{F}$ is an $\mathcal{O}_X$-module of finite type, and
\item the set of weakly associated points of $S$ is locally finite in $S$.
\end{enumerate}
Then $U = \{x \in X \mid \mathcal{F}\text{ flat at }x\text{ over }S\}$
is open in $X$ and $\mathcal{F}|_U$ is flat over $S$ and locally
finitely presented relative to $S$ (see
More on Morphisms, Definition
\ref{more-morphisms-definition-relatively-finitely-presented-sheaf}).
\end{lemma}

\begin{proof}
The question is local on $X$ and $S$. Thus we may assume $X$ and $S$ are affine.
Then we may choose a closed immersion $i : X \to \mathbf{A}^n_S$.
We apply Theorem \ref{theorem-finite-type-flat} to $X' = \mathbf{A}^n_S \to S$
and the quasi-coherent module $\mathcal{F}' = i_*\mathcal{F}$ of finite type
and we find that
$$
U' = \{x' \in X' \mid \mathcal{F}'\text{ flat at }x'\text{ over }S\}
$$
is open in $X'$ and that $\mathcal{F}'|_{U'}$ is of finite presentation.
Since $\mathcal{F}'$ restricts to zero on $X' \setminus i(X)$ and
since $\mathcal{F}'_{i(x)} \cong \mathcal{F}_x$ for all $x \in X$ we see that
$$
U' = i(U) \amalg (X' \setminus i(X))
$$
Hence $U = i^{-1}(U')$ is open. Moreover, it is clear that
$\mathcal{F}'|_{U'} = (i|_U)_*(\mathcal{F}|_U)$.
Hence we conclude that $\mathcal{F}|_U$ is finitely presented
relative to $S$ by More on Morphisms, Lemmas
\ref{more-morphisms-lemma-relative-finite-presentation} and
\ref{more-morphisms-lemma-finite-morphism-relative-finite-presentation}.
\end{proof}

\end{document}
